\documentclass{article}
\usepackage{amsmath,amssymb,amsthm,algorithm,algorithmic}
\usepackage{enumitem}
\usepackage{hyperref}
\newtheorem{assumption}{Assumption}
\newtheorem{theorem}{Theorem}
\newtheorem{lemma}{Lemma}
\newcommand{\EE}{\mathbb{E}}
\newcommand{\RR}{\mathbb{R}}
\newcommand{\norm}[1]{\left\|#1\right\|}
\begin{document}

\section*{Algorithm Name}
\textbf{Federated Averaging with Local SGD (FA‑L‑SGD)}  

The method performs $E$ local stochastic‑gradient steps on each of $K$ clients
between two global aggregations.  The analysis explicitly retains the
\emph{client‑drift} term 
\[
\delta_{k}^{t,i}\;:=\;\nabla F_k\!\bigl(w_{k}^{t,i}\bigr)-\nabla F\!\bigl(w^{t}\bigr),
\]
where $F_k$ is the local objective on client~$k$ and $F=\frac1K\sum_{k=1}^K F_k$ is the global objective.

---------------------------------------------------------------------

\section*{Mathematical Formulation}
\begin{itemize}[leftmargin=*]
  \item \textbf{Global model.} $w^{t}\in\RR^{d}$ at communication round $t=0,1,\dots,T-1$.
  \item \textbf{Local model on client $k$.} $w_{k}^{t,0}=w^{t}$ and for $i=0,\dots,E-1$
  \[
      w_{k}^{t,i+1}=w_{k}^{t,i}-\eta\,g_{k}^{t,i},
      \qquad
      g_{k}^{t,i}:=\nabla F_k\!\bigl(w_{k}^{t,i};\xi_{k}^{t,i}\bigr),
  \]
  where $\xi_{k}^{t,i}$ denotes the random data sample (or mini‑batch) used on
  client $k$ at local step $i$.
  \item \textbf{Global aggregation.}
  After $E$ local steps each client sends $w_{k}^{t,E}$ to the server and the
  server updates
  \[
      w^{t+1}= \frac1K\sum_{k=1}^{K} w_{k}^{t,E}.
  \]
\end{itemize}
All hyper‑parameters are:
\begin{itemize}[leftmargin=*]
  \item $K$ – number of participating clients,
  \item $E$ – number of local SGD steps per round,
  \item $\eta>0$ – learning rate (possibly constant or decreasing).
\end{itemize}

---------------------------------------------------------------------

\section*{Pseudocode}
\begin{algorithm}[H]
\caption{FA‑L‑SGD}
\begin{algorithmic}[1]
\STATE \textbf{Input:} $w^{0}$, learning rate $\eta$, local steps $E$, rounds $T$.
\FOR{$t=0$ \TO $T-1$}
    \FOR{each client $k=1,\dots,K$ \textbf{in parallel}}
        \STATE $w_{k}^{t,0}\gets w^{t}$
        \FOR{$i=0$ \TO $E-1$}
            \STATE Sample mini‑batch $\xi_{k}^{t,i}$
            \STATE $g_{k}^{t,i}\gets \nabla F_k\!\bigl(w_{k}^{t,i};\xi_{k}^{t,i}\bigr)$
            \STATE $w_{k}^{t,i+1}\gets w_{k}^{t,i}-\eta\,g_{k}^{t,i}$
        \ENDFOR
        \STATE Send $w_{k}^{t,E}$ to server
    \ENDFOR
    \STATE $w^{t+1}\gets \frac1K\sum_{k=1}^{K} w_{k}^{t,E}$ \COMMENT{global averaging}
\ENDFOR
\STATE \textbf{Output:} $\displaystyle \bar w_T=\frac1T\sum_{t=0}^{T-1} w^{t}$
\end{algorithmic}
\end{algorithm}

---------------------------------------------------------------------

\section*{Dry Run Example}
Consider the scalar quadratic $f(w)=(w-3)^{2}$, thus $\nabla f(w)=2(w-3)$.
Take a single client ($K=1$), $E=2$ local steps, learning rate $\eta=0.1$,
and initialise $w^{0}=0$.

\begin{center}
\begin{tabular}{c|c|c|c}
Step & $w$ before update & $g=\nabla f(w)$ & $w$ after update \\
\hline
$t=0,i=0$ & $0$ & $2(0-3)=-6$ & $0-0.1(-6)=0.6$ \\
$t=0,i=1$ & $0.6$ & $2(0.6-3)=-4.8$ & $0.6-0.1(-4.8)=1.08$ \\
\hline
Aggregation & $w^{1}=w_{1}^{0,2}=1.08$ & $f(1.08)=(1.08-3)^{2}=3.6864$ &  \\
\hline
$t=1,i=0$ & $1.08$ & $2(1.08-3)=-3.84$ & $1.08-0.1(-3.84)=1.464$ \\
$t=1,i=1$ & $1.464$ & $2(1.464-3)=-3.072$ & $1.464-0.1(-3.072)=1.7712$ \\
\hline
Aggregation & $w^{2}=1.7712$ & $f(1.7712)=1.492$ &
\end{tabular}
\end{center}
The objective value decreases from $9$ to $3.69$ and then to $1.49$, illustrating
the effect of local SGD steps before averaging.

---------------------------------------------------------------------

\section*{Assumptions}
\begin{assumption}[Smoothness]\label{as:smooth}
Each local loss $F_k$ is $L$–smooth:
\[
\forall w,v\in\RR^{d},\qquad
F_k(v)\le F_k(w)+\langle\nabla F_k(w),v-w\rangle+\frac{L}{2}\norm{v-w}^{2}.
\]
Consequently $F$ is also $L$–smooth.
\end{assumption}

\begin{assumption}[Unbiased Stochastic Gradients]\label{as:unbiased}
For every $k,t,i$,
\[
\EE_{\xi_{k}^{t,i}}\!\bigl[g_{k}^{t,i}\mid w_{k}^{t,i}\bigr]=\nabla F_k\!\bigl(w_{k}^{t,i}\bigr).
\]
\end{assumption}

\begin{assumption}[Bounded Variance]\label{as:var}
There exists $\sigma^{2}>0$ s.t.
\[
\EE_{\xi_{k}^{t,i}}\!\bigl[\norm{g_{k}^{t,i}-\nabla F_k(w_{k}^{t,i})}^{2}\mid w_{k}^{t,i}\bigr]
\le\sigma^{2},\qquad\forall k,t,i.
\]
\end{assumption}

\begin{assumption}[Client Drift Bounded]\label{as:drift}
Define the drift $\delta_{k}^{t,i}:=
\nabla F_k(w_{k}^{t,i})-\nabla F(w^{t})$.
There exists $G_{\delta}>0$ s.t.
\[
\EE\!\bigl[\norm{\delta_{k}^{t,i}}^{2}\bigr]\le G_{\delta}^{2},
\qquad\forall k,t,i.
\]
\end{assumption}

\begin{assumption}[Learning Rate]\label{as:lr}
The stepsize satisfies $0<\eta\le \frac{1}{4L\max\{1,E\}}$.
\end{assumption}

---------------------------------------------------------------------

\section*{Main Convergence Theorem}
\begin{theorem}[Non‑convex convergence of FA‑L‑SGD]\label{thm:main}
Let Assumptions \ref{as:smooth}–\ref{as:lr} hold.  Denote
\[
\bar w_T:=\frac1T\sum_{t=0}^{T-1} w^{t}.
\]
Then for any $T\ge1$,
\[
\frac1T\sum_{t=0}^{T-1}\EE\!\bigl[\norm{\nabla F(w^{t})}^{2}\bigr]
\;\le\;
\underbrace{\frac{2\bigl(F(w^{0})-F^{\star}\bigr)}{\eta T}}_{\text{optimization term}}
\;+\;
\underbrace{4L\eta^{2}E^{2}\sigma^{2}}_{\text{stochastic term}}
\;+\;
\underbrace{4L\eta^{2}E^{2}G_{\delta}^{2}}_{\text{drift term}}.
\]
Consequently, choosing a constant stepsize
\[
\eta = \Theta\!\Bigl(\frac{1}{\sqrt{T}}\Bigr)
\]
yields
\[
\frac1T\sum_{t=0}^{T-1}\EE\!\bigl[\norm{\nabla F(w^{t})}^{2}\bigr]
= O\!\Bigl(\frac{1}{\sqrt{T}}\Bigr).
\]
\end{theorem}

---------------------------------------------------------------------

\section*{Theoretical Properties}
\begin{enumerate}
\item \textbf{Stability.}  The bound requires $\eta\le \frac{1}{4L\max\{1,E\}}$,
ensuring that the local SGD iterates remain in a region where the smoothness
inequality is tight.
\item \textbf{Hyper‑parameter sensitivity.}
    \begin{itemize}
        \item Larger $E$ reduces communication but inflates the drift and
        stochastic terms as $E^{2}$ (cf. Theorem \ref{thm:main}).
        \item The variance $\sigma^{2}$ and drift bound $G_{\delta}^{2}$ appear
        only multiplicatively with $\eta^{2}E^{2}$, so a smaller stepsize
        mitigates their impact.
    \end{itemize}
\item \textbf{Comparison to baselines.}
    \begin{itemize}
        \item With $E=1$ the bound collapses to the standard SGD rate
        $O(1/\sqrt{T})$.
        \item When $G_{\delta}=0$ (identical local objectives) the drift term
        disappears and the rate matches that of centralized minibatch SGD with
        batch size $K$.
    \end{itemize}
\end{enumerate}

\section*{Discussion}
The theorem shows that, despite non‑convexity and heterogeneous data,
FA‑L‑SGD attains the same $O(1/\sqrt{T})$ decay of the average squared gradient
as vanilla SGD, at the price of an additional $E^{2}$ factor multiplying the
variance and drift constants.  This factor quantifies the trade‑off between
communication efficiency (large $E$) and statistical efficiency (small $E$).

---------------------------------------------------------------------

\section*{Preliminaries (Definitions and Lemmas)}
\begin{lemma}[Smoothness inequality]\label{lem:smooth}
If $F$ is $L$‑smooth then for any $u,v\in\RR^{d}$
\[
F(v)\le F(u)+\langle\nabla F(u),v-u\rangle+\frac{L}{2}\norm{v-u}^{2}.
\]
\end{lemma}
\begin{lemma}[Descent lemma for one local step]\label{lem:local}
For any client $k$, round $t$, and local step $i$,
\[
\EE\!\bigl[F_k(w_{k}^{t,i+1})\mid w_{k}^{t,i}\bigr]
\le
F_k(w_{k}^{t,i})-\eta\bigl\|\nabla F_k(w_{k}^{t,i})\bigr\|^{2}
+\frac{L\eta^{2}}{2}\bigl(\sigma^{2}+\|\nabla F_k(w_{k}^{t,i})\|^{2}\bigr).
\]
\end{lemma}
\begin{proof}
Apply Lemma~\ref{lem:smooth} with $u=w_{k}^{t,i}$, $v=w_{k}^{t,i+1}=w_{k}^{t,i}-\eta g_{k}^{t,i}$,
take expectation conditioned on $w_{k}^{t,i}$ and use Assumptions
\ref{as:unbiased}–\ref{as:var}.  Rearranging yields the claimed inequality. \qedhere
\end{proof}

---------------------------------------------------------------------

\section*{Main Proof (Step‑by‑Step Derivation)}
We prove Theorem~\ref{thm:main}.  Throughout the proof expectations are taken
with respect to all randomness up to the indicated iteration.

\noindent\textbf{Notation.}
\[
\bar g^{t}:=\frac1K\sum_{k=1}^{K} g_{k}^{t,E-1},
\qquad
\Delta^{t}:=w^{t+1}-w^{t}.
\]

\noindent\textbf{Step 1 (One‑round descent of the global objective).}
\begin{align}
\EE\!\bigl[F(w^{t+1})\mid w^{t}\bigr]
&\overset{(a)}{\le}
\EE\!\Bigl[F(w^{t})+\big\langle\nabla F(w^{t}),\Delta^{t}\big\rangle
+\frac{L}{2}\norm{\Delta^{t}}^{2}\,\Big|\,w^{t}\Bigr] \label{eq:step1}\\
&=F(w^{t})+\big\langle\nabla F(w^{t}),\EE[\Delta^{t}\mid w^{t}]\big\rangle
+\frac{L}{2}\EE\!\bigl[\norm{\Delta^{t}}^{2}\mid w^{t}\bigr].
\end{align}
\begin{description}
\item[Explanation of (a):] Apply Lemma~\ref{lem:smooth} with $u=w^{t}$,
$v=w^{t+1}=w^{t}+\Delta^{t}$ and take conditional expectation. \hfill[Step 1]
\end{description}

\noindent\textbf{Step 2 (Expression for $\Delta^{t}$).}
From the aggregation rule,
\[
\Delta^{t}= \frac1K\sum_{k=1}^{K}\bigl(w_{k}^{t,E}-w^{t}\bigr)
= \frac1K\sum_{k=1}^{K}\sum_{i=0}^{E-1}\bigl(w_{k}^{t,i+1}-w_{k}^{t,i}\bigr)
= -\eta\frac1K\sum_{k=1}^{K}\sum_{i=0}^{E-1} g_{k}^{t,i}.
\tag{Step 2}
\]

\noindent\textbf{Step 3 (Conditional expectation of $\Delta^{t}$).}
Using unbiasedness (Assumption~\ref{as:unbiased}) and the definition of the
drift $\delta_{k}^{t,i}$,
\begin{align}
\EE[\Delta^{t}\mid w^{t}]
&= -\eta\frac1K\sum_{k=1}^{K}\sum_{i=0}^{E-1}
\EE\!\bigl[g_{k}^{t,i}\mid w^{t}\bigr] \nonumber\\
&= -\eta\frac1K\sum_{k=1}^{K}\sum_{i=0}^{E-1}
\bigl(\nabla F_k(w_{k}^{t,i})\bigr) \nonumber\\
&= -\eta\sum_{i=0}^{E-1}\Bigl(\nabla F(w^{t})+\underbrace{\frac1K\sum_{k=1}^{K}
\delta_{k}^{t,i}}_{\text{average drift}}\Bigr) \nonumber\\
&= -\eta E\,\nabla F(w^{t})
-\eta\sum_{i=0}^{E-1}\frac1K\sum_{k=1}^{K}\delta_{k}^{t,i}.
\tag{Step 3}
\end{align}

\noindent\textbf{Step 4 (Plugging Step 3 into the descent inequality).}
Insert \eqref{eq:step1} and the expression for $\EE[\Delta^{t}\mid w^{t}]$:
\begin{align}
\EE\!\bigl[F(w^{t+1})\mid w^{t}\bigr]
&\le F(w^{t})+\big\langle\nabla F(w^{t}),-\eta E\nabla F(w^{t})\big\rangle
   -\eta\Big\langle\nabla F(w^{t}),\sum_{i=0}^{E-1}\frac1K\sum_{k}\delta_{k}^{t,i}\Big\rangle
   +\frac{L}{2}\EE\!\bigl[\norm{\Delta^{t}}^{2}\mid w^{t}\bigr] \nonumber\\
&= F(w^{t})-\eta E\|\nabla F(w^{t})\|^{2}
   -\eta\Big\langle\nabla F(w^{t}),\sum_{i=0}^{E-1}\frac1K\sum_{k}\delta_{k}^{t,i}\Big\rangle
   +\frac{L}{2}\EE\!\bigl[\norm{\Delta^{t}}^{2}\mid w^{t}\bigr]. \tag{Step 4}
\end{align}

\noindent\textbf{Step 5 (Bounding the cross term).}
Using Cauchy–Schwarz and Assumption~\ref{as:drift},
\begin{align}
\Big|\big\langle\nabla F(w^{t}),\sum_{i=0}^{E-1}\frac1K\sum_{k}\delta_{k}^{t,i}\big\rangle\Big|
&\le \|\nabla F(w^{t})\|\;
      \Big\|\sum_{i=0}^{E-1}\frac1K\sum_{k}\delta_{k}^{t,i}\Big\| \nonumber\\
&\le \|\nabla F(w^{t})\|\;
      \sum_{i=0}^{E-1}\frac1K\sum_{k}\|\delta_{k}^{t,i}\|
 \le \|\nabla F(w^{t})\|\;E\,G_{\delta}. \tag{Step 5}
\end{align}
Taking expectation and applying $ab\le\frac{a^{2}}{2}+\frac{b^{2}}{2}$,
\[
\EE\!\Bigl[\eta\big\langle\nabla F(w^{t}),\sum_{i=0}^{E-1}\frac1K\sum_{k}\delta_{k}^{t,i}\big\rangle\Bigr]
\le \frac{\eta E}{2}\EE\!\bigl[\|\nabla F(w^{t})\|^{2}\bigr]
   +\frac{\eta E}{2}G_{\delta}^{2}. \tag{Step 5'}
\]

\noindent\textbf{Step 6 (Bounding $\EE[\|\Delta^{t}\|^{2}]$).}
From Step 2,
\[
\Delta^{t}= -\eta\sum_{i=0}^{E-1}\frac1K\sum_{k=1}^{K} g_{k}^{t,i}.
\]
Hence,
\begin{align}
\EE\!\bigl[\norm{\Delta^{t}}^{2}\mid w^{t}\bigr]
&= \eta^{2}\EE\!\Bigl[\Big\|\sum_{i=0}^{E-1}\frac1K\sum_{k} g_{k}^{t,i}\Big\|^{2}\Big| w^{t}\Bigr] \nonumber\\
&\le \eta^{2}E\sum_{i=0}^{E-1}\EE\!\Bigl[\Big\|\frac1K\sum_{k} g_{k}^{t,i}\Big\|^{2}\Big| w^{t}\Bigr] \quad (\text{Jensen}) \nonumber\\
&\le \eta^{2}E\sum_{i=0}^{E-1}
\Bigl(\underbrace{\Big\|\frac1K\sum_{k}\nabla F_k(w_{k}^{t,i})\Big\|^{2}}_{\le 2\|\nabla F(w^{t})\|^{2}+2G_{\delta}^{2}}
    +\frac{\sigma^{2}}{K}\Bigr) \nonumber\\
&\le \eta^{2}E^{2}\bigl(2\|\nabla F(w^{t})\|^{2}+2G_{\delta}^{2}\bigr)
   +\eta^{2}E^{2}\frac{\sigma^{2}}{K}. \tag{Step 6}
\end{align}
The inequality $\big\|\frac1K\sum_k\nabla F_k\big\|^{2}
\le 2\|\nabla F(w^{t})\|^{2}+2G_{\delta}^{2}$ follows from
$\nabla F_k(w_{k}^{t,i})=\nabla F(w^{t})+\delta_{k}^{t,i}$ and
$(a+b)^{2}\le 2a^{2}+2b^{2}$.

\noindent\textbf{Step 7 (Putting everything together).}
Insert the bounds from Steps 5' and 6 into the descent inequality of
Step 4, then take total expectation:
\begin{align}
\EE\!\bigl[F(w^{t+1})\bigr]
&\le \EE\!\bigl[F(w^{t})\bigr]
   -\eta E\EE\!\bigl[\|\nabla F(w^{t})\|^{2}\bigr]
   +\frac{\eta E}{2}\EE\!\bigl[\|\nabla F(w^{t})\|^{2}\bigr]
   +\frac{\eta E}{2}G_{\delta}^{2} \nonumber\\
&\qquad
   +\frac{L}{2}\Bigl(\eta^{2}E^{2}2\EE\!\bigl[\|\nabla F(w^{t})\|^{2}\bigr]
                     +\eta^{2}E^{2}2G_{\delta}^{2}
                     +\eta^{2}E^{2}\frac{\sigma^{2}}{K}\Bigr) \nonumber\\
&= \EE\!\bigl[F(w^{t})\bigr]
   -\underbrace{\Bigl(\eta E-\frac{\eta E}{2}-L\eta^{2}E^{2}\Bigr)}_{\displaystyle
        \alpha}\EE\!\bigl[\|\nabla F(w^{t})\|^{2}\bigr]
   +\underbrace{\Bigl(\frac{\eta E}{2}+L\eta^{2}E^{2}\Bigr)}_{\displaystyle \beta}
      G_{\delta}^{2}
   +\frac{L\eta^{2}E^{2}\sigma^{2}}{2K}. \tag{Step 7}
\end{align}
By Assumption~\ref{as:lr}, $\eta\le\frac{1}{4L\max\{1,E\}}$ implies
$\alpha\ge \frac{\eta E}{2}$ and $\beta\le \eta E$.
Thus,
\[
\EE\!\bigl[F(w^{t+1})\bigr]
\le \EE\!\bigl[F(w^{t})\bigr]
   -\frac{\eta E}{2}\EE\!\bigl[\|\nabla F(w^{t})\|^{2}\bigr]
   +\eta E\,G_{\delta}^{2}
   +\frac{L\eta^{2}E^{2}\sigma^{2}}{2K}. \tag{Step 7'}
\]

\noindent\textbf{Step 8 (Summation over $t$).}
Sum the inequality in Step 7' for $t=0,\dots,T-1$ and telescope:
\begin{align}
\frac{\eta E}{2}\sum_{t=0}^{T-1}\EE\!\bigl[\|\nabla F(w^{t})\|^{2}\bigr]
&\le F(w^{0})-F^{\star}
   +T\eta E\,G_{\delta}^{2}
   +\frac{L\eta^{2}E^{2}\sigma^{2}}{2K}T. \tag{Step 8}
\end{align}
Dividing by $T\eta E/2$ yields
\[
\frac1T\sum_{t=0}^{T-1}\EE\!\bigl[\|\nabla F(w^{t})\|^{2}\bigr]
\le \frac{2\bigl(F(w^{0})-F^{\star}\bigr)}{\eta ET}
   +2G_{\delta}^{2}
   +\frac{L\eta E\sigma^{2}}{K}. \tag{Step 8'}
\]

\noindent\textbf{Step 9 (Final simplification).}
Since $K\ge1$ and $L\le 2L$, we can absorb constants to obtain the cleaner
statement of Theorem~\ref{thm:main}:
\[
\frac1T\sum_{t=0}^{T-1}\EE\!\bigl[\|\nabla F(w^{t})\|^{2}\bigr]
\le
\frac{2\bigl(F(w^{0})-F^{\star}\bigr)}{\eta T}
+4L\eta^{2}E^{2}\sigma^{2}
+4L\eta^{2}E^{2}G_{\delta}^{2}.
\]
\hfill[Step 9]

Choosing $\eta = \frac{c}{\sqrt{T}}$ with a constant $c$ that respects
Assumption~\ref{as:lr} gives the $O(1/\sqrt{T})$ rate claimed in the theorem.
\qed

---------------------------------------------------------------------

\section*{Rate Derivation (With Constants)}
From the bound in Theorem~\ref{thm:main},
let $\eta = \min\Bigl\{\frac{1}{4L\max\{1,E\}},\;\frac{\sqrt{2(F(w^{0})-F^{\star})}}{L^{1/2}E\sqrt{T}}\Bigr\}$.
Then
\[
\frac1T\sum_{t=0}^{T-1}\EE\!\bigl[\|\nabla F(w^{t})\|^{2}\bigr]
\le
\underbrace{O\!\Bigl(\frac{1}{\sqrt{T}}\Bigr)}_{\text{optimization term}}
\;+\;
\underbrace{O\!\Bigl(\frac{E^{2}}{T}\Bigr)}_{\text{variance term}}
\;+\;
\underbrace{O\!\Bigl(\frac{E^{2}}{T}\Bigr)}_{\text{drift term}}.
\]
Thus the dominant term is $O(1/\sqrt{T})$ provided $E=O(1)$ or $E$ grows at most
as $T^{1/4}$.

---------------------------------------------------------------------

\section*{Special Cases}
\begin{enumerate}
\item \textbf{Identical local data ($G_{\delta}=0$).}
    The drift term disappears, and the bound matches centralized mini‑batch
    SGD with batch size $K\!E$.
\item \textbf{Single local step ($E=1$).}
    The factor $E^{2}$ becomes $1$, recovering the classic non‑convex SGD
    rate $O(1/\sqrt{T})$.
\item \textbf{Full participation, infinite local steps ($E\to\infty$).}
    The bound degrades as $E^{2}$; in practice one caps $E$ to keep the
    $E^{2}$ term comparable to the optimisation term.
\end{enumerate}

\end{document}